\begin{lemma}[The exact terminal route under final-tree stability]
Let $G$ be a finite loopless multigraph with individual edge identities,
and let $\mathcal E$ be an extended Tashkinov tree with distinguished
uncolored edge $e_0$ and positive rung number $n$. Write $T_n$ for its
terminal ladder tree, $F$ for its series' final output, $\varphi_n$ for
its terminal coloring, and $D_n$ for its computed terminal stored set.
Let $\sigma$ be proper, leave exactly $e_0$ uncolored, and be
$(F,D_n,\varphi_n)$-stable. Then $\varphi_n$ is proper and leaves exactly
$e_0$ uncolored, $T_n$ is an exact initial segment of $F$ with strictly
smaller length, and $\sigma$ is $(T_n,D_n,\varphi_n)$-stable.
Moreover, precisely the following terminal-route alternative holds:
either the terminal kind is PE and $F$ is a $\sigma$-closure of $T_n$,
or the terminal kind is SE or RE and there are an exact edge $f$ and
an ordered seed $S$ with terminal connecting-edge set exactly $\{f\}$,
$S=T_n+f$ in the exact one-edge-adjunction sense, and $F$ a
$\sigma$-closure of $S$.
\end{lemma}
\begin{proof}
The positivity field of \noderef{ExtendedTashkinovTree} makes $q=n-1$
a valid noninitial-step index. Inspect the certificate at index $n$ in
\noderef{TashkinovSeries}, using \noderef{TashkinovSeriesStep}. The initial
constructor has index zero and is impossible here. In SE the current
coloring is the selected coloring; its properness and exactly-one-uncolored
properties are fields of \noderef{MaximumDefectiveChoice}. In PE these
properties are explicit current-coloring fields. In RE the current
coloring equals the preceding coloring and inherits both preceding-coloring
fields. Thus $\varphi_n$ has both properties in all three cases.
Every case explicitly records $T_n\preccurlyeq F$ and
$\operatorname{length}(T_n)<\operatorname{length}(F)$.
The PE case records a $\varphi_n$-closure from $T_n$ to $F$.
The SE case records the singleton connecting set $\{f\}$ for the selected
exact edge and an actual ordered tree $S$ satisfying
\noderef{TreeSequenceAdjoinEdge}$(T_n,f,S)$, followed by a
$\varphi_n$-closure of $S$ to $F$. The RE case records these same data
with the exact edge from its revisiting certificate. No edge is replaced
by its endpoint pair in this argument.

We transport each of these recorded closures. Fix its source $S$ and
expand \noderef{TashkinovClosure} into an augmentation route and closedness
of $F$. By \noderef{TashkinovAugmentationSequence}, the route is a finite
chain. For a head step $A\to B$ followed by its suffix to $F$,
\noderef{TashkinovAugmentationSequenceInitialSegment} gives exact initial
segments $A\preccurlyeq F$ and $B\preccurlyeq F$.
\noderef{TreeSequenceInitialSegmentVertexSetSubset} puts their vertices
in $V(F)$. The step in \noderef{TashkinovAugmentation} supplies its exact
added edge $g$, its color $c=\varphi_n(g)$, and $c$ missing at some
$v\in V(A)$, by \noderef{TreeMissingColors}. Hence $c$ is missing on
$F$. The attachment law of \noderef{TreeSequence} gives an endpoint of
$g$ in $V(B)\subseteq V(F)$, so $g$ is incident with $F$ by
\noderef{Incident} and \noderef{TreeIncidentEdge}. The edge clause of
\noderef{StableColoring} gives $\sigma(g)=\varphi_n(g)=c$. Its vertex
clause at $v$ says that $c$ is still missing at the same vertex under
$\sigma$. Reuse the unchanged length and exact-prefix clauses of the
step: it is now a $\sigma$-augmentation. Head induction transports the
whole route, including the reflexive route when $S=F$.

To check closedness under $\sigma$, suppose that $g$ is a boundary edge
of $F$ with $\sigma(g)=c$, where $c$ is missing under $\sigma$ at a vertex
$v\in V(F)$. The endpoint definition \noderef{BoundaryEdges} supplies
an endpoint in $V(F)$, so $g$ is incident with $F$.
The vertex clause of stability makes $c$ missing under $\varphi_n$ at
$v$, hence tree-missing on $F$. Both colorings are proper and leave
exactly $e_0$ uncolored, as proved above. Apply the two-sided edge
criterion \noderef{sc1} to the given final-tree stability: its alternative
using the new value $\sigma(g)=c$ yields $\sigma(g)=\varphi_n(g)$.
This contradicts original closedness, in the sense of
\noderef{ColoredVertexSetClosed}, at $g,v,c$. The transported route and
this closedness prove the identical recorded closure under $\sigma$.
In the SE/RE cases we retained the same $f$, singleton connecting set,
and exact ordered seed $S$ throughout.

Finally restrict stability to $T_n$. The exact initial-segment relation
already proved gives $V(T_n)\subseteq V(F)$ by
\noderef{TreeSequenceInitialSegmentVertexSetSubset}. Thus the vertex
clause restricts directly, and an edge incident with $T_n$ is also
incident with $F$. The second inclusion of
\noderef{HkuMissingStoredMonotone}, at the valid index $q=n-1$, states
\[
 \overline\varphi_n(T_n)\cup D_n
 \subseteq\overline\varphi_n(F)\cup D_n.
\]
Consequently every edge value protected by stability on $T_n$ is already
protected by the given stability on $F$. The edge clause therefore
restricts too. These are the two clauses of
\noderef{StableColoring}$(T_n,D_n,\varphi_n,\sigma)$, completing all
claims.
\end{proof}
